% Generated by inst/sup_info/table_and_figure_supplement.R; do not edit by hand.
\begin{singlespace}
\small
\begin{verbatim}
> library(twocpcont)
> ss <- twocpcont_ss(delta1 = 1.8, delta2 = 3.1, sd1 = 9, sd2 = 12,
+                    rho = 0.5, r = 1, alpha = 0.025, beta = 0.13,
+                    method = "univariate")
> ss
  Sample size calculation: two co-primary continuous endpoints
  -------------------------------------------------------------- 
  [Design]
  Mean difference (d1, d2) : 1.800, 3.100
  Std. deviation (s1, s2)  : 9, 12
  Correlation (rho)        : 0.5000
  Allocation ratio (r)     : 1
  Significance level       : 0.02500
  Target power             : 0.8700
  -------------------------------------------------------------- 
  [Computation]
  Method                   : Noniterative closed-form (proposed)
  -------------------------------------------------------------- 
  [Results]
  Sample size (n1, n2)     : 490, 490
  Total sample size (N)    : 980
  -------------------------------------------------------------- 
> ss$n2_cont
[1] 489.664
> twocpcont_power(n1 = ss$n1, n2 = ss$n2, delta1 = 1.8, delta2 = 3.1,
+                 sd1 = 9, sd2 = 12, rho = 0.5, alpha = 0.025)
  Power calculation: two co-primary continuous endpoints
  ---------------------------------------------------------- 
  [Design]
  Sample size (n1, n2)     : 490, 490
  Mean difference (d1, d2) : 1.800, 3.100
  Std. deviation (s1, s2)  : 9, 12
  Correlation (rho)        : 0.5000
  Significance level       : 0.02500
  ---------------------------------------------------------- 
  [Computation]
  Method                   : Bivariate normal (conventional)
  ---------------------------------------------------------- 
  [Results]
  Marginal power (ep. 1)   : 0.8791
  Marginal power (ep. 2)   : 0.9814
  Co-primary power         : 0.8702
  ---------------------------------------------------------- 
> sapply(c(0, 0.3, 0.5, 0.8), function(rho) {
+   twocpcont_ss(delta1 = 1.8, delta2 = 3.1, sd1 = 9, sd2 = 12,
+                rho = rho, alpha = 0.025, beta = 0.13,
+                method = "univariate")$N
+ })
[1] 1000  990  980  960
> kappa <- (3.1 / 12) / (1.8 / 9)
> r_max(kappa, alpha = 0.025, beta = 0.13)
[1] 0.04618346
> kappa_star(epsilon = 0.05, alpha = 0.025, beta = 0.13)
[1] 1.279183
\end{verbatim}
\end{singlespace}
